\section{Where the rubrics do not settle it}

Four questions in this book's territory are genuinely open on the text. Two
more are commonly treated as open and are in fact closed by evidence the
Missal itself supplies. Both kinds are set out here, because a reference that
hides its uncertainties is worse than useless at the moment it is most needed.

\subsection{Open: the funeral Mass on an ordinary Sunday}

\begin{contested}{whether \rgm{406} bars the funeral Mass on a second-class
Sunday}
\begin{rulebox}{\rgm{406}}
Missa exsequialis prohibetur: a) diebus sub nn.~1, 2, 3, 4, 5 et 6 in tabella
praecedentiae recensitis; b) diebus festis de praecepto, inter festa sub n.~11
in tabella praecedentiae comprehensis; c) in anniversario Dedicationis et in
festo Tituli ecclesiae, in qua funus peragitur; d) in festo Patroni principalis
oppidi vel civitatis; e) in festo Tituli et Sancti Fundatoris Ordinis seu
Congregationis, ad quam pertinet ecclesia in qua funus peragitur.
\end{rulebox}

\textbf{The question.} Rows 1--6 catch Christmas, Easter and Pentecost, the
Triduum, Epiphany, the Ascension, Trinity, Corpus Christi, the Sacred Heart,
Christ the King, the Immaculate Conception, the Assumption, the vigil and
octave day of Christmas, and every first-class Sunday. Second-class Sundays are
row 15. Are they caught by clause (b)?

\textbf{Reading one: they are not.} \latin{Comprehensis} is ablative plural
agreeing with \latin{diebus festis de praecepto}, so the clause reads ``on days
of precept \emph{comprised among} the feasts under n.~11'' --- a restriction,
not an addition. Row 11 is ``first-class feasts of the universal Church not
named above,'' which contains no Sunday. On this reading the prohibition adds
nothing but the universal first-class feasts that a given territory keeps as
days of obligation, and an ordinary green Sunday is not barred.

\rgm{407} corroborates it, and does so in a way that is hard to explain
otherwise: ``si vero solemnitas externa alicuius festi fit in dominica, Missa
exsequialis prohibetur die quo fit solemnitas externa, non autem die festo.''
That sentence needs a Sunday on which the funeral Mass would otherwise be
possible; if every Sunday were already barred by \rgm{406}\,b, the clause would
have nothing to forbid. External solemnities \latin{ipso iure} fall on ordinary
Sundays --- Our Lady of the Rosary on the first Sunday of October is the plain
case (\rgm{358}\,b) --- so the clause is about second-class Sundays.

\textbf{Reading two: they are.} \latin{Dies festus de praecepto} was not a term
the rubrics coined, and in contemporaneous Roman usage it included Sundays. The
Sacred Congregation of the Council, by decree of 3 December 1960 --- five months
after the code and effective the same 1 January 1961 --- issued a taxative list
of \latin{festa de praecepto} which opens ``Dominicae I et II Classis'' before
naming the ten first-class feasts of the universal calendar. A drafter using
that phrase in 1960 could reasonably be read as importing that list, with the
row-11 clause identifying which \emph{feasts} qualify rather than excluding
Sundays. On this reading a funeral Mass is barred on every Sunday.

\textbf{Our assessment, and its limits.} The grammar and \rgm{407} both favour
reading one, and we think a careful reader of the Latin will reach it. We do
not think that settles what a parish should do. \rgm{393}\,b independently bars
every Mass of the dead in a church having only one Mass whenever a conventual
burden presses that no other priest can discharge, and the practical obligation
of the parish Mass on Sunday will usually leave no room. The result is that the
textual answer and the pastoral answer diverge, and the Ordinary and the Ordo,
not this book, govern the second.

Two things this argument does \emph{not} establish. The 1960 decree defines
\latin{festa de praecepto} for the obligation of applying the Mass
\latin{pro populo} under cann.~339 \S1 and 466 \S1 of the 1917 Code; it is not
a list of the days on which the faithful are bound to hear Mass, which is a
separate question of canon law that this study does not enter. And nothing
here bears on whether a funeral may be held on a Sunday at all --- only on
which Mass may be said if it is.
\end{contested}

\subsection{Open: how far the sung-Mass restriction reaches}

\begin{contested}{what counts as a \latin{Missa in cantu non conventualis}}
\rg{108} and \rg{111}\,a between them delete every ordinary commemoration from
a sung Mass that is not the conventual Mass, and \rgm{434}\,a repeats the rule
in terms of orations. \rgm{271} defines a Mass as sung ``si sacerdos celebrans
partes ab ipso iuxta rubricas cantandas revera cantu profert'' --- if the
celebrant actually sings the parts the rubrics assign him --- and makes the
\latin{Missa cantata}, sung without sacred ministers, a species of sung Mass.
So a parish \latin{Missa cantata} is caught.

What is not settled is the boundary case that parishes actually meet: a Mass
at which a choir sings the Ordinary and the propers while the celebrant reads
his parts. \rgm{271}'s test is what the \emph{celebrant} does, and on the
text such a Mass is a read Mass with singing, so the ordinary commemorations
stay. But the same rubric's purpose --- and the practice of describing such a
Mass as a sung Mass --- pulls the other way, and a celebrant who begins reading
and then sings the Preface has moved the Mass across the line mid-way.

We take the literal test to be right: the class of the Mass follows what the
celebrant sings, and a Mass at which he sings nothing the rubrics assign him is
a read Mass however much the choir sings. We record that the question is
practically live and that a wrong answer changes the number of collects.

The neighbouring rule has no such ambiguity and is worth stating beside it:
\rgm{457}\,d bars the prayer imposed by the Ordinary ``in Missis in cantu''
without qualification --- every sung Mass, conventual or not.
\end{contested}

\subsection{Open: two questions already argued above}

Two further open questions are argued at the place where a reader meets them,
and are only listed here so that the inventory is complete.

\begin{itemize}
\item \textbf{Whether the Preface of the Trinity reaches a Mass that is not of
the Sunday} --- \rgm{494}\,b's ``in dominicis II classis'' against the
neighbouring rubrics' ``in Missis quae celebrantur eodem tempore.'' Section~7.

\item \textbf{The year with only twenty-three Sundays after Pentecost} ---
\rg{18} does not expressly legislate that shortfall. The calendar computation
profile requires a fail-closed result pending a competent dated Ordo.
\end{itemize}

\subsection{Closed by evidence: the Purification is a feast of the Lord}

\begin{contested}{resolved --- whether 2 February takes the place of a
second-class Sunday}
\rg{16}\,a gives a feast of the Lord of the first or second class occurring on
a second-class Sunday the Sunday's own place ``cum omnibus iuribus et
privilegiis,'' with no commemoration of the Sunday. Whether the Purification of
the Blessed Virgin Mary, a second-class feast, counts as a \latin{festum
Domini} for that rule is a natural question: the feast is titled of Our Lady,
and \rg{120}\,b assigns it white as a Marian feast.

The Missal answers it in its own rubric, printed under the title on 2 February:
``\latin{Festum Purificationis beatae Mariae Virginis habetur tamquam festum
Domini}.'' The code corroborates it twice from the Preface rules: \rgm{484}\,a
gives the Purification the Preface of the Nativity ``tamquam propria,'' and
\rgm{495} excepts it by name from the Preface of the Blessed Virgin. So when 2
February falls on a Sunday after Epiphany or a Sunday after Pentecost, the
Purification takes the day, no commemoration of the Sunday is made, and the
Creed is nevertheless said, because \rgm{475}\,a attaches the Creed to the
Sunday even when its Office has yielded.

The related permission is separate and should not be confused with this. Where
the liturgical action proper to 2 February --- the blessing and procession of
candles --- is by apostolic approval transferred to a Sunday, an external
solemnity belongs to the feast \latin{ipso iure}, but only for the Mass which
follows that blessing and procession (\rgm{358}\,c). An external solemnity is
not a transfer of the feast, and the Office stays where it was.
\end{contested}

\subsection{Closed by evidence: the Missal is not self-sufficient}

\begin{contested}{resolved --- where the Gloria rule actually lives}
\rgm{431}\,a is the primary rule for the Gloria: it is said ``in Missis quae
respondent Officio diei, quotiescumque ad Matutinum dictus est hymnus Te
Deum.'' The rule for the Te Deum is \rgb{237}--\rgb{238}. Those numbers are in
the second part of the code, \work{Rubricae generales Breviarii Romani}, which
n.~5 of \work{Rubricarum instructum} directed to be printed at the front of the
Breviary and not of the Missal. The 1962 Vatican typical edition follows that
instruction exactly: it prints the heading ``II.\ --- \textsc{Rubricae
generales Breviarii Romani} (nn.~138--268)'' and, beneath it, the two words
``\latin{Hic omittuntur}.''

The Missal therefore states a rule whose antecedent it does not contain, and
nothing in the Missal supplies it. This is not an editorial slip: it is what
the motu proprio ordered, on the assumption that the celebrant has said Matins
and knows whether he said the Te Deum. But it means that a server, a layman
with a hand missal, or anyone reasoning from the Missal alone cannot apply
\rgm{431}\,a, and it is the reason the summary of \rgb{237}--\rgb{238} is set out
at the head of Section~8. The same limitation affects \rgm{424}--\rgm{425},
\rgm{299} and every other Missal rubric that turns on the shape of the Office.
\end{contested}
